\section{Agreement of neighboring linear labels}\label{sec:neighbor-labels}

The affine profiles of Proposition~\ref{prop:affine-profiles} may initially
have different linear parts, or \emph{labels}, on different sheets.
We now exclude that possibility
for sheets joined by one generator step of the enlarged chart width.
The argument uses only the local affine profiles.  In particular, it does
not use a global sublinear estimate, or equality of their affine offsets.

We retain the polytope $\mathcal C$, the domain $\mathcal U_\delta$, and
the profiles $f_{r,x,n}$ from Definition~\ref{def:chart-profiles}.
All limits below are along sequences $r\to\infty$.  A factor $e^{o(r)}$
has one common vanishing envelope throughout each finite construction.
The constants defining the family of Lipschitz quasi-isometry pairs are
fixed.

\begin{proposition}[Neighboring labels]\label{prop:neighbor-labels}
There is $\delta_1>0$, depending only on the fixed group and the uniform
quasi-isometry bounds, with the following property.  Fix
$0<\delta\leq\delta_1$.  Let $F$, $x$, and the following suffixes and steps
vary with $r$:
\[
 \begin{gathered}
 a\sim\frac{1_{B_r(\mathcal C)}}{\vol B_r(\mathcal C)}\,dn,
 \qquad n=ak_0,\qquad
 k_0\in B_r((1+\delta)\mathcal C,e^{o(r)}),\\
 k=\exp(tX_i),
 \quad |t|\leq e^{(1+\delta)r s_i(\mathcal C)+o(r)}.
 \end{gathered}
\]
The suffix and step are deterministic.  Every pair in the support of a
limiting joint law of $(f_{r,x,n},f_{r,x,nk})$ has the form
\[
 (A\pi z+b_1,\ A\pi z+b_2)\qquad(z\in\mathcal U_\delta)
\]
for a common $A\in\mathcal A$ and some $b_1,b_2\in W$.
\end{proposition}

The contradiction comes from reflecting an inward height segment at a
facet of $\mathcal C$.  If the proposition failed for arbitrarily small
enlargements, the two limiting affine maps would agree on that facet,
and their relative linear map would be a reflection.  An inward source
segment $E$ would then have two opposite target images relative to the
common value at its facet endpoint.  Each image has volume exponent
$P(E)$, while their union has exponent $2P(E)$.

We will alternate finite sequences of switches on the two sheets.
Their source terminal points remain reachable over $E$, so their
separated count has exponent at most $P(E)$.  To force the larger
target exponent, we retain a large joint volume of successive endpoints.
The polynomial image estimate then turns the corresponding target
increments into terminal volume with exponent $2P(E)$.  The next lemma
supplies the intervening passage from separated endpoint tuples to
volume in lifted $N$ coordinates, allowing $N\cap D\ne\{1\}$.
Write $q:N\rtimes D\to G$ for the multiplication map.

\begin{lemma}[Volume from endpoint tuples]\label{lem:tuple-lift-volume}
Fix a positive integer $h$ and compact sets $E_1,\ldots,E_h\subset W$
containing $0$.  Suppose that, for each $r$, there are $M_r$ tuples
$(y_1,\ldots,y_h)\in G^h$, separated by a sufficiently large fixed
constant in the maximum product metric.  Let $y_0$ be a common basepoint.
Suppose each tuple admits lifts relative to $y_0$ of the form
\[
 y_0^{-1}y_j=q(U_1\cdots U_j,d_j),\qquad
 U_j\in B_r(E_j,e^{o(r)}),\qquad
 |d_j|_D\leq Cr,\quad \|\pi d_j\|=o(r),
\]
with common bounds for all tuples.  Then there is a measurable set
$\mathcal T_r\subset N^h$ of increments satisfying the same box bounds,
with possibly larger subexponential factors, such that
\begin{equation}\label{eq:tuple-lift-volume}
 \vol_{N^h}(\mathcal T_r)\geq M_r e^{-o(r)}.
\end{equation}
Every product $u_1\cdots u_h$ from $\mathcal T_r$ has an accompanying
$d\in D$ with $|d|_D=O(r)$ and $\pi d=o(r)$ whose projection lies a
bounded distance from one of the listed terminal points $y_0^{-1}y_h$.
Moreover, the $N$-coordinate volume of all such pairs projecting into any
one ball of fixed radius in $G$ is at most $e^{o(r)}$.
\end{lemma}

\begin{proof}
Choose one lift for each endpoint of each tuple.  Thicken each total lift
$(U_1\cdots U_j,d_j)$ by multiplying it on the right by a fixed small
compact neighborhood in $N\rtimes D$.  Its projected diameter is uniformly
bounded, because $q$ is a homomorphism.  The resulting product
neighborhoods in $(N\rtimes D)^h$ are disjoint: their projections are
contained in disjoint neighborhoods of the separated tuples.  Each has
the same positive Haar volume.  Here product Haar on $N\rtimes D$ is
both left and right invariant, by unimodularity and the determinant-one
$D$ action on $N$.

All their $D$ coordinates lie in a ball of radius $O(r)$, of polynomial
Haar volume.  Fubini therefore shows that their projection to the total
$N$ coordinates has volume at least $M_r$ divided by a polynomial in
$r$.  The map
\[
 (v_1,\ldots,v_h)\longmapsto
 (v_1,v_1^{-1}v_2,\ldots,v_{h-1}^{-1}v_h)
\]
preserves product Haar: successively, each changed variable undergoes a
left translation.  Apply it to this projected set to obtain
$\mathcal T_r$.

To check its bounds, write a thickened total $N$ coordinate as
$(U_1\cdots U_j)\eta_j$.  Right multiplication of its full lift by a
bounded factor makes
$\eta_j=\Ad_{d_j}(a_j)$ for $a_j$ in a fixed compact subset of $N$.
The small height and polynomial unipotent bound give
$\|\log\eta_j\|\leq e^{o(r)}$.  The new increment is
$\eta_{j-1}^{-1}U_j\eta_j$, with $\eta_0=1$.
Since $0\in E_j$, all word widths of $E_j$ are nonnegative, and the box
product estimates put this new increment in $B_r(E_j,e^{o(r)})$.
Every projected total coordinate has at least one accompanying $D$
tuple from the thickened set; this proves the assertion about terminal
products as well as \eqref{eq:tuple-lift-volume}.

For completeness, consider two admissible pairs $(u,d)$ and
$(\widetilde u,\widetilde d)$ whose projections lie in the same ball of
fixed radius.  Lift their bounded group difference as $q(a,b)$ with
$a\in N$ and $b\in D$ in fixed compact sets.  Equality in $G$ gives
\begin{equation}\label{eq:tuple-overlap}
 u^{-1}\widetilde u=\Ad_d(a)z,
 \qquad z=d b\widetilde d^{-1}\in D\cap N.
\end{equation}
The first factor has log norm at most $e^{o(r)}$.  The second has
$D$-length $O(r)$ and hence polynomial log coordinates, by
Lemma~\ref{lem:d-scaling}.  The logarithm of their product is therefore
bounded by $e^{o(r)}$.  Once one admissible pair is fixed, all the other
$N$ coordinates in the ball lie in a left translate of an $N$ log ball
of radius $e^{o(r)}$.  Its Haar volume is $e^{o(r)}$.  This is also the
specific instance of Lemma~\ref{lem:lift-multiplicity} needed here.
Notice that \eqref{eq:tuple-overlap} bounds the entire overlap, without
assuming it is central or discrete.
\end{proof}

\begin{proof}[Proof of Proposition~\ref{prop:neighbor-labels}]
Both profiles are covered by Proposition~\ref{prop:affine-profiles},
since multiplying the allowed suffix by $k$ preserves its enlarged box
bound.  Suppose the asserted equality of labels fails for arbitrarily
small positive $\delta$.  We explain first the limiting geometry and
then compare source and target volumes.

\paragraph{A reflection at a facet.}
Choose violating $\delta\to0$.  Passing to a subsequence fixes the
letter $X_i$ and distinct labels $A_1,A_2\in\mathcal A$.  Compactness
of the profiles allows their offsets to converge to $b_1,b_2$.
Diagonalize in the scales so that there are Borel sets
\[
 Q_r=B_r(\mathcal C)k_0,\qquad S_r\subset Q_r,\qquad
 \frac{\vol S_r}{\vol Q_r}\geq e^{-o(r)},
\]
on which the two profiles converge to
$A_1\pi z+b_1$ and $A_2\pi z+b_2$ on every required compact set over
$\mathcal C$.  To justify the mass assertion, fix each approximation
tolerance first.  A neighborhood of a point in a limiting support has
positive probability.  Take its scale sufficiently large that the
negative logarithm of this probability divided by $r$ is arbitrarily
small, then decrease the tolerance.  A compact exhaustion permits any
later fixed finite collection of paths.  The domains
$\mathcal U_\delta$ exhaust the inverse image of
$\operatorname{int}\mathcal C$; equicontinuity on compact subsets of
$D_\infty$ extends this convergence to the boundary along the same
diagonal.
Also $k_0,k\in B_r(\mathcal C,e^{o(r)})$ on this diagonal.

At a height in the facet
\[
 H_i=\{w\in\mathcal C:\alpha_i(w)=s_i(\mathcal C)\},
\]
the points on the two sheets are a distance $o(r)$ apart: conjugate the
step by a $D$ position at that height and use
Lemma~\ref{lem:log-distortion}.  Thus the limiting affine maps agree on
the relative interior of the entire facet.  Their linear parts satisfy
\[
 A_1^{-1}A_2 v=v\qquad(v\in\ker\alpha_i).
\]
Choose an inner product invariant under the finite group $\mathcal A$.
An orthogonal map fixing a hyperplane pointwise is either the identity
or its orthogonal reflection.  Since $A_1\ne A_2$, the map
$A_1^{-1}A_2$ is this reflection.

Let $p_0$ be the contact point in the relative interior of $H_i$ and
let $v$ be its reflection normal, chosen with $\alpha_i(v)>0$.
For a sufficiently small fixed $\epsilon>0$ put
\[
 E_0=[-\epsilon v,0],\qquad E=p_0+E_0\subset\mathcal C,
 \qquad E'_1=A_1E_0,\qquad E'_2=A_2E_0=-A_1E_0.
\]
The common affine value at $p_0$ is the target pivot.  Translation
invariance and segment homogeneity of $P$ give
\begin{equation}\label{eq:reflection-doubling}
 P(E'_1)=P(E'_2)=P(E)=\epsilon p(v)>0,
 \qquad P(E'_1\cup E'_2)=2P(E).
\end{equation}
Figure~\ref{fig:reflection} depicts these source and target segments.

\begin{figure}[tbp]
\centering
\begin{tikzpicture}[>=Stealth,font=\small]
 \begin{scope}
  \node[font=\small\bfseries] at (1.65,2.5) {Source heights};
  \fill[blue!5] (-.1,-.15) rectangle (2.75,2.12);
  \draw[gray,dashed] (2.75,-.2)--(2.75,2.15);
  \node[anchor=east] at (2.65,1.78) {$\alpha_i=s_i(\mathcal C)$};
  \draw[blue!65!black,very thick] (.55,.95)--(2.75,.95);
  \fill (.55,.95) circle (1.7pt);
  \fill (2.75,.95) circle (1.7pt);
  \node[below] at (.55,.95) {$p_0-\varepsilon v$};
  \node[right] at (2.75,.95) {$p_0$};
  \node[above] at (1.55,.95) {$E=p_0+E_0$};
  \node at (1.6,-.55) {$P(E)=P(E_0)>0$};
 \end{scope}
 \begin{scope}[xshift=7cm]
  \node[font=\small\bfseries] at (0,2.5) {Target heights relative to the pivot};
  \draw[blue!65!black,very thick,->] (0,.95)--(-2.05,.95);
  \draw[orange!80!black,very thick,->] (0,.95)--(2.05,.95);
  \fill (0,.95) circle (1.7pt);
  \node[below] at (0,.95) {$0$};
  \node[above] at (-1.1,.95) {$E'_1=A_1E_0$};
  \node[above] at (1.1,.95) {$E'_2=-A_1E_0$};
  \node at (0,-.55) {$P(E'_1\cup E'_2)=2P(E_0)$};
 \end{scope}
\end{tikzpicture}
\caption{The height geometry of the reflection contradiction. The two affine profiles agree at the facet pivot and send the inward segment to opposite target segments. Alternating bundles detect the volume exponent of their union. The picture is schematic: the exponential counts come from bracket-generated boxes, and the volume argument accounts for the nonuniqueness of lifts to $N\rtimes D$.}
\label{fig:reflection}
\end{figure}


\paragraph{Source bundles with a large joint volume.}
At a pivot $d_p\in D$ of height $rp_0$ and $D$-length $O(r)$, alternate
two sorts of bundles.  A first-sheet bundle updates the pair state by
$n\mapsto nu$.  A second-sheet bundle updates it by
$n\mapsto nkuk^{-1}$; on that sheet this is the generator movement
$nk\mapsto nku$.  Each bundle consists of a fixed finite sequence of
generator increments with widths $e^{rs_j(E)}$ and with independent
difference-of-interval parameters before restrictions are imposed.

We choose the source schemes and their number before restricting
parameters to good sheets.  Pass to a subsequence on which the normalized
bracket laws for $E$, $E'_1$, $E'_2$, and $E'_1\cup E'_2$ converge,
together with the bounded linear embeddings from either individual
target normalization into the union normalization.

The number of bundles is determined by multiplication in the target.
Consider the polynomial map multiplying alternating increments from
the two target boxes, with output normalized to $E'_1\cup E'_2$.
For a sufficiently large fixed number $h$ of increments, its limiting
map is submersive.  Indeed, every union-scaled generator attains its
support maximum in one of the two sets.  Its limiting vector therefore
belongs to the image of the corresponding limiting embedding.  These
vectors bracket-generate the union-normalized limiting algebra.
Finite products of exponentials from the two embedding images have
a submersive endpoint map: their conjugates span the generated Lie
algebra, and one inserts variable increments with those conjugates as
differentials.  Inserting identity factors makes the product alternating.
Fix such an $h$.  This rank property belongs to the full multiplication
map, before any target increments have been selected.

For each of the $h$ source bundles,
Lemma~\ref{lem:product-schemes} supplies an $E$-normalized generator
scheme with a submersive limiting product map.  Represent its scalar
parameters as independent differences of interval parameters.  The
schemes and their total number of switches are now fixed.  These
choices depend only on the normalized algebra data, not on $S_r$ or
the support tolerances.  Passing to this subsequence preserves the
profile approximations and the bound $\vol S_r/\vol Q_r\geq e^{-o(r)}$.
All the scaled $D$-legs lie in one compact subset over $\mathcal C$,
including their facet pivots.  The boundary convergence just established
therefore controls their profiles.

Apply backward recurrence to the individual switches in these bundles.
For a first-sheet switch with interval law $\sigma$, the increment law
is $\sigma*\check\sigma$.  For a second-sheet switch it is
$\sigma_k*\check\sigma_k$, where $\sigma_k=(\Ad_k)_*\sigma$ is the conjugated
interval law.
Thus all switches have the form required by
Lemma~\ref{lem:backward-recurrence}, with a depth independent of $r$.

Every recurrence denominator in Lemma~\ref{lem:backward-recurrence}
is at most $e^{o(r)}\vol Q_r$.  Indeed the supports in question are
products of a fixed number of $k_0^{\pm1}$, $k^{\pm1}$, and generator
intervals fitting the $\mathcal C$ widths, all with subexponential
factors.  The box calculus bounds the volume of the corresponding
support enlargements by $e^{o(r)}\vol B_r(\mathcal C)$.
Hence backward recurrence supplies a nonempty Borel subset of $S_r$
such that each initial state in that subset has conditional successful
parameter probability at least $e^{-o(r)}$, with every intermediate
pair state in $S_r$.  Fix one such state $n_0$.

Apply Corollary~\ref{cor:restricted-products} to the map from all bundle
parameters to the tuple of their products \emph{before} the conjugations
by $k$ and before accumulation.  The normalized map is the product of
$h$ submersive limiting polynomial maps.  The normalized pairs of
interval parameters have a fixed bounded joint density with respect
to Lebesgue measure.  Thus the successful probability just proved
gives successful Lebesgue measure at least $e^{-o(r)}$, even though
the restricted coordinates are correlated.  The image consequently
contains a compact set of $N^h$ Haar volume at least
$e^{rhP(E)-o(r)}$.  Conjugation by $k$ preserves Haar on $N$, and the
change from successive increments to successive states is triangular
and Haar preserving.  Consequently the admissible tuples
$(n_1,\ldots,n_h)$ of bundle-end states have the same lower volume.
This application takes place before conjugation, so it requires no
coefficient bound for a conjugation written in $E$-normalized
coordinates.

By Lemma~\ref{lem:fixed-sheet-packing}, the corresponding tuples
\[
 \bigl(F(xn_1d_p),\ldots,F(xn_hd_p)\bigr)
\]
have a separated-count lower bound $e^{rhP(E)-o(r)}$.
Take the source separation sufficiently large that the quasi-isometry
preserves a fixed target separation.  On the source side each terminal
point $xn_hd_p$ is reachable from $xn_0d_p$ by a path of length $O(r)$
whose relative heights lie in $r(E_0+o(1))$.  Use the $D$ legs over $E$
and return to the pivot after each generator switch.  Second-sheet
bundles use $o(r)$ bridges at the pivot.  Lemma~\ref{lem:path-packing}
therefore bounds the covering number of all source terminal points by
\begin{equation}\label{eq:source-reflection-count}
 e^{rP(E)+o(r)}
\end{equation}
at every fixed radius.

\paragraph{Target products have twice the terminal exponent.}
Put $y_j=F(xn_jd_p)$.  The successful profiles imply that the transition
from $y_{j-1}$ to $y_j$ has a path of length $O(r)$ with relative height
set $r(E'_{\tau(j)}+o(1))$, where $\tau(j)$ alternates between $1$ and
$2$.  All pivot height differences are $o(r)$, since the two limiting
affine profiles agree at $p_0$ and every intermediate state is good.
Choose lifts of these paths.  On accumulating them, the preceding
$D$ factors have height $o(r)$ and length $O(r)$; their conjugations
preserve the indicated box widths up to subexponential factors.
The endpoint lifts thus have the form
\[
 (u'_1\cdots u'_j,d'_j),\qquad
 u'_j\in B_r(E'_{\tau(j)},e^{o(r)}),\qquad
 |d'_j|_D=O(r),\quad \pi d'_j=o(r).
\]
Apply Lemma~\ref{lem:tuple-lift-volume} to the separated target tuples.
It gives a measurable increment set of Haar volume at least
$e^{rhP(E)-o(r)}$.  Dividing each factor by its box determinant shows
that this set has Lebesgue measure at least $e^{-o(r)}$ in the joint
individual-box log coordinates.  These coordinates lie in a cube of
radius $e^{o(r)}$.

Apply the alternating multiplication map used to choose $h$, with
output normalized to $E'_1\cup E'_2$.  Its degree is fixed, its
coefficients are bounded by the box calculus, and its limit is
submersive by that earlier choice.  These properties concern the full
map and therefore remain available for the measurable increment set
just obtained.

Lemma~\ref{lem:polynomial-image}, applied to the measurable increment
set, gives a compact set of terminal $N$ coordinates of volume at least
\[
 e^{rP(E'_1\cup E'_2)-o(r)}=e^{2rP(E)-o(r)}.
\]
Every such terminal coordinate has an accompanying $D$ coordinate of
length $O(r)$ and height $o(r)$ projecting near an image terminal.
The last assertion of Lemma~\ref{lem:tuple-lift-volume} bounds its
volume over any fixed-radius projected ball by $e^{o(r)}$.
But \eqref{eq:source-reflection-count} and the quasi-isometry cover all
image terminals by only $e^{rP(E)+o(r)}$ fixed-radius balls.  This is a
contradiction because $P(E)>0$.  The labels therefore agree for every
sufficiently small fixed $\delta$, as asserted.
\end{proof}

The preceding comparison leaves the offsets undetermined.  For fixed
$\delta$, a generator step can require a height exceeding the facet
by order $\delta r$; co-location alone would therefore leave an error
of order $\delta$ in the normalized affine offsets.  The next section
eliminates this error by a separate quotient-volume comparison.
